
De que forma o processo contábil impacta a tomada de decisão dos usuários da contabilidade?

Há que se considerar que os usuários da contabilidade podem ser muitos e apresentar diferentes pontos de vista.

O processo contábil não se limita à elaboração dos relatórios contábeis. Ele se desenvolve dentro das entidades e envolve pessoas de toda a organização. Se entendermos que os processos das empresas que envolvem fluxos financeiros geram informações que devem ser tratadas no processo contábil, qualquer decisão tomada com base nesses registros será dependente ou estará relacionada a esse processo.

Um colaborador, ao receber uma mercadoria, realizar a conferência das quantidades e registrar a entrada, está, em um sentido amplo, participando de um processo contábil. Assim, qualquer decisão que utilize esses dados é uma decisão relacionada ao processo contábil. A decisão de contratar ou não um funcionário, por exemplo, pode requerer a análise dos dados da folha de pagamento. Por essa abrangência, podemos dizer que quase todas as operações de uma empresa são baseadas em dados que participam do processo contábil ou do processo de contabilização.

Analisando de forma mais restrita os usuários dos demonstrativos contábeis, também é importante distinguir seus diferentes tipos e interesses. Entendo que podemos classificá-los em três grupos: aqueles que não participam da entidade, mas têm interesse nela; os administradores e colaboradores; e os sócios.

Quanto aos usuários da contabilidade que não têm, ou ainda não têm, vínculo com a entidade, podemos citar os investidores do mercado de ações. Estes procuram nos demonstrativos contábeis evidências para decidir sobre a gestão de seus ativos, passando ou não a participar da entidade. Também se enquadram nesse grupo as instituições financeiras e os fornecedores, de modo geral, que necessitam decidir a forma como se relacionarão com as entidades, avaliando, entre outros aspectos, sua capacidade de geração de caixa. Da mesma forma, órgãos de fiscalização do governo podem decidir sobre a conformidade da entidade com a legislação com base em informações contábeis.

Existem também os administradores das entidades, que utilizam os demonstrativos como forma de avaliar se as estratégias e decisões tomadas no passado apresentaram os resultados esperados e, a partir dessa análise, decidir sobre ações e planos futuros. Os gestores também podem avaliar os demonstrativos contábeis e, considerando as discricionariedades permitidas pelas normas contábeis, decidir pela sua reelaboração, de forma que evidenciem melhor os resultados da gestão.

Por fim, existem os sócios, que podem utilizar os demonstrativos contábeis para medir o desempenho da entidade e decidir sobre a continuidade de sua participação societária ou, dependendo de sua capacidade de influência, decidir pela substituição da diretoria.
